\documentclass[a4paper]{article}
%\usepackage{exercise}
%um nur aufgaben zu zeigen
\usepackage[noanswer]{exercise} 
\usepackage{../images/preamble}
\usepackage{rotating}
\usetikzlibrary{decorations.pathmorphing}
\usetikzlibrary{decorations.markings}
\usetikzlibrary{arrows}
\usetikzlibrary{shapes.geometric}
\newcommand{\midarrow}{\tikz \draw[-triangle 90] (0,0) -- +(.02,0);}
\usepackage{xcolor}
%\usepackage{draftwatermark}
%\SetWatermarkText{\textsc{Draft 2}}
%\SetWatermarkScale{3}
%\SetWatermarkColor{red!30}

\usepackage[printwatermark]{xwatermark}
%\newsavebox\mybox
%\savebox\mybox{\tikz[color=red,opacity=0.3]\node{\textsc{Entwurf}};}
%\newwatermark*[
%allpages,
%angle=45,
%scale=10,
%xpos=-4cm,
%ypos=4cm
%]{\usebox\mybox}
\pagestyle{fancy}
\fancyhead[L]{\includegraphics[width=2cm]{../images/logo_scaled.pdf}}
\fancyhead[R]{\textsc{Lösung Aufgabenserie 5}}


\begin{document}
	\vspace*{-1cm}
	\parbox{4cm}{\vspace{-0.2cm}\includegraphics[width=5cm]{../images/logo_scaled.pdf}}
	\parbox{10.6cm}{\setstretch{2.0} \centering{ \huge \textsf{Aufgabenserie 5 
			}}\\
			Abgabe: 30. November \\ \vspace*{-.5cm} }
		\vspace{0.5cm}

\thispagestyle{empty}
\begin{framed}
	\noindent
	\scriptsize
	Die Aufgaben sollten bis zum \textbf{30. Novermber} bearbeitet werden. Die Lösungen schickt ihr an \href{mailto:physikrolf@gmail.com}{physikrolf@gmail.com}.
	Jede Aufgabe hat eine bestimmte Anzahl an erreichbaren Punkten. Wie viele das sind, müsst ihr raten. Versucht, die Lösungen so genau wie möglich aufzuschreiben. Für besonders schnelle/gute/witzige Lösungen kann es Bonuspunkte geben.\\ Die aktuellen Aufgaben sowie alle alten Aufgabenserien mit Lösungen findet ihr auch auf \url{pankratius.github.io/rolf}. %\\\textit{Zu jeder Aufgabe gibt es jetzt Tipps. Die sollten beim Lösen der Aufgaben helfen.\\ Sollte das so sein macht bitte in euren Lösungen kenntlich, dass bestimmte Schritte von den Tipps und nicht von euch kommen. Darauf gibt es keinen Abzug, es ist nur für uns gut zu wissen.}
\end{framed}

\noindent

\begin{Exercise}[label = voltmeter, difficulty = 3, origin = EstPhO 2003, title = Viele Voltmeter]
	In der abgebildeten Schaltung sind 15 identische Voltmeter und 15 unterschiedliche Ampermeter verbaut, und an eine Batterie angeschlossen. \\
	Das erste Voltmeter zeigt eine Spannung von $U_1 = 9~\mathrm{V}$ an, das erste Ampermeter einen Strom von $I_1 = 2.9~\mathrm{mA}$ und das zweite einen Strom von $I_2 = 2.6~\mathrm{mA}$. Wie groß ist die Summe der Spannungen, die die anderen Voltmeter anzeigen?
\end{Exercise}
\begin{figure}[h]
	\centering
	\tikzset{component/.style={draw,thick,circle,fill=white,minimum size =0.75cm,inner sep=0pt}}
	\begin{tikzpicture}
		\draw (0,1) to[battery1] (0,-1);
		\draw (0,1) to node[component]{$I_1$} (2,1);
		\draw (2,1) to node[component]{$U_1$}(2,-1);
		\draw (2,-1) -- (0,-1);
		\draw (2,1) to node[component]{$I_2$} (4,1);
		\draw (4,1) to node[component]{$V_2$} (4,-1);
		\draw (4,-1) -- (2,-1);
		\draw (4,1) -- (5,1);
		\draw (4,-1) -- (5,-1);
		\node at (5.5,1) {...};
		\node at (5.5,-1) {...};
		
		\draw (6,1) to node[component]{$I_{14}$} (8,1);
		\draw (8,1) to node[component]{$U_{14}$} (8,-1);
		\draw (8,-1) -- (6,-1);
		\draw (8,1) to node[component]{$I_{15}$} (10,1);
		\draw (10,1) to node[component]{$U_{15}$} (10,-1);
		\draw (10,-1) -- (8,-1);
		
		\foreach \n in {2,4,8}{
			\filldraw (\n,1) circle (2pt);
			\filldraw (\n,-1) circle (2pt);
			}
		
		
	\end{tikzpicture}
\end{figure}
\begin{Exercise}[title = Schwimmende Glasschale, origin = {Lucas Rettenmeier, Orpheus e.V.}, difficulty = 3, label = glas]
	Eine dünnwandige, ideal wärmeleitende  zylindrische Glasschale der Höhe $H = 15~\mathrm{cm}$ schwimmt mit dem Boden nach unten auf einer Wasseroberfläche, und taucht dabei bis zur Hälfte in das Wasser ein.\\
	Bestimme, wie tief die Glasschale eintaucht, wenn man sie umdreht, und ins Wasser tut. Berechne außerdem, wie tief man sie ins Wasser tauchen muss, damit sie nicht mehr schwimmt, sondern untergeht.\\
	Die Wasserdichte ist $\rho = 1000~\mathrm{\nicefrac{kg}{m^3}}$ und der äußere Luftdruck beträgt $p_0 = 101 325~\mathrm{Pa}$. \\ Die eingeschlossene Luft kann man näherungsweise als ein sog. \textit{ideales Gas} betrachten. Für ein solches gilt in diesem Fall die \textit{ideale Gasgleichung}, die besagt, dass das Verhältnis
			$\frac{pV}{T}$
	die ganze Zeit konstant bleibt. Dabei ist $p$ der Druck der Luft in dem Gefäß, $V$ das Volumen der eingeschlossenen Luft und $T$ die Temperatur. Gleichzeitig kann angenommen werden, dass sich die Lufttemperatur während des Umdrehens nicht ändert. 
\end{Exercise}
\begin{Answer}[ref = glas]
	Die Wasserschale soll, wenn man sie wie beschrieben mit der Öffnung nach oben ins Wasser stellt, schwimmen.\\
	Also müssen Gewichtskraft $F_g$ und Auftriebskraft $F_a$ gleich sein,
	\begin{equation}\label{glas:auftrieb1}
		m g = \frac{1}{2}Vg\rho,
	\end{equation}
	wobei $m$ die Masse der Schale ist und $V$ das Gesamtvolumen.\\
	Dreht man die Schale jetzt um, dringt Wasser ein. Es muss aber wieder das Gleichgewicht von Gewichtskraft und Auftriebskraft gelten. Dafür ist erneut das verdrängte Volumen relevant. Um das Auszurechnen, brauchen wir die Eindringtiefe der Schale in das Wasser $h$ und die Wasserhöhe im Glas $x$. Dann gilt
	\begin{equation}\label{glas:auftrieb2}
		mg = \frac{V}{H}\left(h-x\right) g \rho.
	\end{equation}
	Gleichsetzen von \eqref{glas:auftrieb1} und \eqref{glas:auftrieb2} führt auf
	\begin{equation}\label{glas:mechheight}
		\frac{1}{2}Vg\rho = \frac{V}{H}\left(h-x\right) g \rho \Rightarrow h-x = \frac{H}{2}
	\end{equation}
	Wir können jetzt die ideale Gasgleichung verwenden, um den weitere Informationen über die Höhe des Wasserstandes im Glas, der durch den Druck im Glas bestimmt wird, zu erhalten.\\
	Dazu können wir annehmen, dass sich die Temperatur beim Umdrehen nicht ändert, weil die Aufgabenstellung sagt, dass alles ideal wärmeleitend ist.\\
	Vor dem Umdrehen betrug das Volumen der  Luft im Glas genau $V$, und der Druck entsprach dem äußeren Luftdruck $p_0$.\\
	Nach dem Umdrehen beträgt das Volumen der Luft im Glas genau $V \cdot \frac{H-h}{H}$, da das der Teil ist, der nicht mit Wasser gefüllt ist. Der Druck entspricht der Summe des äußeren Luftdrucks $p_0$ und des Schweredrucks durch die Wassersäule der Höhe $h-x$, also $p_0 + \rho g \left(h-x\right)$. Setzt man alles in die ideale Gasgleichung ein, dann kommt man auf
	\begin{equation}
		\frac{p_0 V}{T} = \frac{\left(p_0 + \rho g \left(h-x\right) \right) V \cdot \frac{H-h}{H}}{T}.
	\end{equation}
	Wir können jetzt das Ergebnis aus \eqref{glas:mechheight} für $h-x$ einsetzen, und die entstehende Gleichung dann nach der neuen Eintauchtiefe $h$ lösen
	\begin{equation}\boxed{
		p_0= \left(p_0 + \frac{\rho g H}{2}\right) \frac{H-h}{H} \Rightarrow h = H \left(\frac{3}{2} - \frac{2p_0}{2p_0 + \rho g H}\right) \approx 7.6~\mathrm{cm}.}
	\end{equation}
	Damit die Schale nicht mehr auftaucht, sondern bis zum Boden sinkt, muss das enthaltene Luftvolumen bis auf einen Wert kleiner als $\frac{V}{2}$ komprimiert werden. \\
	Es ergibt sich also für diese Tiefe
	\begin{equation*}
	\boxed{
		\frac{H}{2}+ \frac{p_0}{\rho g} \approx 10~\mathrm{m}
		}
	\end{equation*}
\end{Answer}
\begin{Exercise}[difficulty = 3, origin = {IV. IPhO, 1970}, label = sphmi, title = Teleskop]
	In einem Teleskop wird ein sphärischer Spiegel mit einem Krümmungsradius von $R = 2~\mathrm{m}$ verwendet. Im Brennpunkt des Spiegels befindet sich der Strahlungsempfänger $E$ in Form einer runden Scheibe.  Diese liegt senkrecht zur optischen Achse des Teleskops.\\
	Welche Größe muss der Empfänger haben, damit er die gesamt Strahlung empfängt, die vom Spiegel mit dem Radius von $r = 0,25~\mathrm{m}$ reflektiert wird?
\end{Exercise}


\end{document}